\documentclass[11pt,a4paper]{article}

% ── Codificación y fuentes ────────────────────────────────────────────────────
\usepackage[utf8]{inputenc}
\usepackage[T1]{fontenc}
\usepackage{lmodern}
\usepackage[spanish]{babel}

% ── Geometría y espaciado ─────────────────────────────────────────────────────
\usepackage[top=2.5cm, bottom=2.5cm, left=2.8cm, right=2.8cm]{geometry}
\usepackage{setspace}
\setstretch{1.15}
\usepackage{parskip}

% ── Tablas ────────────────────────────────────────────────────────────────────
\usepackage{booktabs}
\usepackage{tabularx}
\usepackage{array}
\usepackage{longtable}
\usepackage{enumitem}

% ── Colores y cajas ───────────────────────────────────────────────────────────
\usepackage[dvipsnames]{xcolor}
\usepackage{tcolorbox}
\tcbuselibrary{skins, breakable}

% ── Cabeceras y pies de página ────────────────────────────────────────────────
\usepackage{fancyhdr}
\usepackage{lastpage}
\pagestyle{fancy}
\fancyhf{}
\fancyhead[L]{\small\textbf{Informe de Evaluación} --- Física para Bioanalistas}
\fancyhead[R]{\small ybramaris maribao 36150392}
\fancyfoot[C]{\small Página \thepage\ de \pageref{LastPage}}
\renewcommand{\headrulewidth}{0.4pt}

% ── Hiperenlaces ──────────────────────────────────────────────────────────────
\usepackage[colorlinks=true, linkcolor=NavyBlue, urlcolor=NavyBlue]{hyperref}

% ── Paleta de colores personalizados ─────────────────────────────────────────
\definecolor{desempeno_excelente}{HTML}{1A6B2F}
\definecolor{desempeno_bueno}{HTML}{2E5A9C}
\definecolor{desempeno_suficiente}{HTML}{8B6914}
\definecolor{desempeno_deficiente}{HTML}{C0392B}
\definecolor{desempeno_insuficiente}{HTML}{7B241C}
\definecolor{grisFondo}{HTML}{F5F5F5}
\definecolor{azulTitulo}{HTML}{1B2A6B}
\definecolor{verdeFortaleza}{HTML}{1A6B2F}
\definecolor{naranjaMejora}{HTML}{A04000}

% ── Estilos de cajas ─────────────────────────────────────────────────────────
\tcbset{
  cajaTitulo/.style={
    enhanced, breakable,
    colback=azulTitulo!8, colframe=azulTitulo,
    fonttitle=\bfseries\large, coltitle=white,
    attach boxed title to top left={yshift=-2mm, xshift=4mm},
    boxed title style={colback=azulTitulo},
    top=6pt, bottom=6pt, left=8pt, right=8pt,
  },
  cajaFortaleza/.style={
    enhanced, breakable,
    colback=verdeFortaleza!6, colframe=verdeFortaleza!60,
    fonttitle=\bfseries, coltitle=verdeFortaleza,
    top=4pt, bottom=4pt, left=8pt, right=8pt,
  },
  cajaMejora/.style={
    enhanced, breakable,
    colback=naranjaMejora!6, colframe=naranjaMejora!60,
    fonttitle=\bfseries, coltitle=naranjaMejora,
    top=4pt, bottom=4pt, left=8pt, right=8pt,
  },
  cajaRecomendacion/.style={
    enhanced, breakable,
    colback=grisFondo, colframe=gray!50,
    fonttitle=\bfseries, coltitle=black,
    top=4pt, bottom=4pt, left=8pt, right=8pt,
  },
}

% ─────────────────────────────────────────────────────────────────────────────
\begin{document}

% ══ PORTADA ══════════════════════════════════════════════════════════════════
\begin{center}
  {\Large\bfseries\color{azulTitulo} Universidad de Oriente/Núcleo Sucre --- Física para Bioanalistas}\\[4pt]
  {\large Informe de Evaluación Preliminar}\\[12pt]
  \rule{\linewidth}{1.2pt}\\[6pt]
  {\LARGE\bfseries ybramaris maribao 36150392}\\[4pt]
  \rule{\linewidth}{0.6pt}
\end{center}

% ══ FICHA TÉCNICA ════════════════════════════════════════════════════════════
\vspace{4pt}
\begin{tcolorbox}[cajaTitulo, title=Datos del informe]
\begin{tabular}{@{}llll@{}}
  \textbf{Estudiante:}  & ybramaris maribao 36150392  &
  \textbf{Fecha:}       & 2026-06-26 \\[4pt]
  \textbf{Archivo PDF:} & \texttt{ybramaris\_maribao\_36150392.pdf}  &
  \textbf{Páginas:}     & 5 \\
\end{tabular}
\end{tcolorbox}

% ══ RESULTADO GLOBAL ═════════════════════════════════════════════════════════
\vspace{6pt}
\begin{center}
  \begin{tcolorbox}[
    enhanced,
    colback=white, colframe=desempeno_bueno,
    width=0.62\linewidth,
    boxrule=2pt,
    halign=center, valign=center,
    top=8pt, bottom=8pt,
  ]
    \centering
    {\large\bfseries Puntaje sugerido}\\[6pt]
    {\Huge\bfseries\color{desempeno_bueno} 7.7/10}\\[6pt]
    {\large Nivel de desempeño: \textbf{\color{desempeno_bueno} Bueno}}
  \end{tcolorbox}
\end{center}

% ══ RESUMEN GENERAL ══════════════════════════════════════════════════════════
\section*{Resumen general}
El estudiante demuestra una buena comprensión de los principios fundamentales de la física de fluidos relevantes para el bioanálisis, aplicando correctamente las fórmulas para presión hidrostática, empuje, continuidad y la ley de Poiseuille. La mayoría de los problemas fueron planteados correctamente con una adecuada extracción de datos y conversión de unidades. Sin embargo, se observan problemas recurrentes con la precisión en los cálculos numéricos, especialmente en aquellos que involucran potencias y notación científica, y la falta de completitud en una de las preguntas.

% ══ FORTALEZAS Y ÁREAS DE MEJORA ════════════════════════════════════════════
\vspace{4pt}
\begin{minipage}[t]{0.48\linewidth}
  \begin{tcolorbox}[cajaFortaleza, title={Fortalezas}]
\begin{itemize}[leftmargin=1.5em, itemsep=2pt]
  \item Correcta identificación y aplicación de los principios físicos y fórmulas relevantes (presión hidrostática, principio de Arquímedes, continuidad, Bernoulli, Poiseuille).
  \item Buena capacidad para extraer datos de los enunciados de los problemas y realizar conversiones de unidades (ej. mm a m).
  \item Presentación clara y organizada de los pasos para la mayoría de los problemas.
  \item Fuerte comprensión conceptual en las preguntas 1, 2, 3 y 5.
\end{itemize}
  \end{tcolorbox}
\end{minipage}
\hfill
\begin{minipage}[t]{0.48\linewidth}
  \begin{tcolorbox}[cajaMejora, title={Areas de mejora}]
\begin{itemize}[leftmargin=1.5em, itemsep=2pt]
  \item Precisión en los cálculos numéricos, especialmente aquellos que involucran exponentes y notación científica.
  \item Completitud en la respuesta a todas las partes de una pregunta (ej. finalizar el análisis proporcional).
  \item Atención al detalle en la escritura de unidades (ej. unidad de densidad en la Pregunta 1).
\end{itemize}
  \end{tcolorbox}
\end{minipage}

% ══ OBSERVACIONES POR PREGUNTA ═══════════════════════════════════════════════
\section*{Observaciones por pregunta}

\begin{longtable}{>{\bfseries}p{2.8cm} p{1.6cm} p{7.0cm} p{2.8cm}}
  \toprule
  \textbf{Pregunta} & \textbf{Puntaje} & \textbf{Evaluación} & \textbf{Errores detectados} \\
  \midrule
  \endhead
  \bottomrule
  \endfoot
    \textbf{Pregunta 1: Presión Hidrostática y Venosa} & 1.8/2.0 & El estudiante aplicó correctamente la fórmula de presión hidrostática para determinar la altura requerida. La extracción de datos y la sustitución en la fórmula fueron adecuadas. El resultado numérico es correcto. & \textit{Error al escribir la unidad de densidad (kg/m$^{2}$ en lugar de kg/m$^{3}$), aunque el cálculo posterior sugiere que se usó el valor correcto para la densidad volumétrica.} \\
    \midrule
    \textbf{Pregunta 2: Principio de Arquímedes} & 2.0/2.0 & Excelente resolución. El estudiante calculó correctamente la fuerza de empuje, el volumen del objeto y su densidad, aplicando de manera impecable el principio de Arquímedes. Todos los pasos, cálculos y unidades son correctos. & \textit{---} \\
    \midrule
    \textbf{Pregunta 3: Hemodinámica (Continuidad y Bernoulli)} & 1.9/2.0 & Muy buena aplicación de las ecuaciones de continuidad y Bernoulli. El cálculo de la velocidad en la zona estrecha fue correcto. La fórmula de Bernoulli se aplicó adecuadamente y el resultado final es correcto. & \textit{Se detectó un error menor en un cálculo intermedio (530 * -1.35 = -415.5, cuando debería ser -715.5). Sin embargo, el resultado final (13284.5 Pa) es consistente con el cálculo correcto de -715.5, lo que sugiere un error de transcripción o un error de calculadora que fue corregido implícitamente en el paso final.} \\
    \midrule
    \textbf{Pregunta 4: Ley de Poiseuille (Análisis Proporcional)} & 1.0/2.0 & El estudiante identificó correctamente la relación de Poiseuille (Q $\propto$ r$^{4}$) y la reducción del radio (r\_final = 0.80 * r\_inicial). Sin embargo, no completó el cálculo para determinar el porcentaje de caudal que se mantendría, dejando la pregunta sin una respuesta numérica final. & \textit{Respuesta incompleta: No se calculó el cambio porcentual en el caudal.} \\
    \midrule
    \textbf{Pregunta 5: Flujo Viscoso y Resistencia} & 1.0/2.0 & El estudiante identificó correctamente la Ley de Poiseuille y la despejó para la diferencia de presión. La sustitución de los datos en la fórmula fue correcta. & \textit{Error significativo en el cálculo del término r$^{4}$ en el denominador ($\pi$ * (4 x 10$^{-}$$^{4}$)$^{4}$). El valor calculado para el denominador es incorrecto por un factor de 100, lo que lleva a un resultado final numéricamente erróneo.} \\
    \midrule
\end{longtable}

% ══ ERRORES DETECTADOS ═══════════════════════════════════════════════════════
\section*{Análisis de errores}

\subsection*{Errores conceptuales}
\begin{itemize}[leftmargin=1.5em, itemsep=2pt]
  \item Ninguno mayor. El error de unidad en la Pregunta 1 (kg/m$^{2}$ para densidad) podría ser un error conceptual si el estudiante realmente lo cree, pero el uso posterior en el cálculo sugiere un error de escritura.
\end{itemize}

\subsection*{Errores procedimentales}
\begin{itemize}[leftmargin=1.5em, itemsep=2pt]
  \item Error al escribir la unidad de densidad en la Pregunta 1 (kg/m$^{2}$ en lugar de kg/m$^{3}$).
  \item Error menor de cálculo en un paso intermedio de la Pregunta 3 (530 * -1.35).
  \item La Pregunta 4 está incompleta: no se realizó el cálculo final del cambio porcentual del caudal.
  \item Error significativo en el cálculo del denominador (término r$^{4}$) en la Pregunta 5, afectando el resultado final.
\end{itemize}

% ══ RECOMENDACIONES AL ESTUDIANTE ════════════════════════════════════════════
\section*{Retroalimentación para el estudiante}
\begin{tcolorbox}[cajaRecomendacion, title={Dirigido a: ybramaris maribao 36150392}]
Estimado estudiante, has demostrado una sólida comprensión de los principios fundamentales de la física de fluidos. Tu capacidad para identificar las fórmulas correctas y configurar los problemas es una gran fortaleza. Para mejorar, te sugiero prestar más atención a la precisión en tus cálculos numéricos, especialmente con exponentes y notación científica. Revisa siempre tus pasos y utiliza la calculadora con cuidado. Además, asegúrate de responder completamente a todas las partes de la pregunta, como en el caso del análisis proporcional, donde el cálculo final era necesario. Practicar más con problemas que involucren potencias y unidades te ayudará a consolidar estos aspectos.
\end{tcolorbox}

% ══ NOTA PARA EL PROFESOR ════════════════════════════════════════════════════
\section*{Nota para el profesor}
\begin{tcolorbox}[
  enhanced, breakable,
  colback=yellow!8, colframe=yellow!60!black,
  fonttitle=\bfseries, coltitle=black,
  title={Observaciones para revisión manual},
  top=4pt, bottom=4pt, left=8pt, right=8pt,
]
El estudiante muestra un buen nivel de comprensión conceptual en la mayoría de los temas. Los errores principales son de índole procedimental y de cálculo, especialmente en las Preguntas 4 y 5. En la Pregunta 3, el error intermedio de cálculo parece haber sido corregido implícitamente para llegar al resultado final correcto, lo que podría indicar un buen entendimiento pero un descuido en la transcripción. La Pregunta 4 está conceptualmente bien planteada pero incompleta. Sugiero enfatizar la importancia de la precisión en los cálculos y la revisión de las respuestas finales.
\end{tcolorbox}

% ══ PIE DE INFORME ════════════════════════════════════════════════════════════
\vspace{12pt}
\begin{center}
  \footnotesize\color{gray}
  Informe generado automáticamente por el sistema de evaluación con IA (gemini-2.5-flash) y soporte LaTex. \\
  Este documento es una evaluación \textbf{preliminar} para ser revisado por el profesor antes de ser entregado al 
  estudiante. \\
  Generado el 2026-06-26 \textbullet\ BIOANALISIS Sistema Académico de Evaluación.
  \end{center}

\end{document}
